\documentstyle[%


righttag%


,11pt%


,amssymb%


,russian%               remove in Eng.


,izvru%                 Set izven% in Eng.


]{amsart}





%





\newcommand{\Span}{\operatorname{span}}


\newcommand{\tts}[1]{}


\numberwithin{equation}{section}


\renewcommand{\baselinestretch}{0.98}





\theoremstyle{plain}


\theorembodyfont{\it}


\newtheorem{athm}{\hskip\parindent Теорема}


\renewcommand{\theathm}{\Alph{athm}}





\theoremstyle{plain}


\theorembodyfont{\it}


\newtheorem{mainthm}{\hskip\parindent Основная теорема}


\renewcommand{\themainthm}{}





\theoremstyle{definition}


\theorembodyfont{\rm}


\newtheorem{notenonum}{\hskip\parindent Замечание}


\renewcommand{\thenotenonum}{}





%\hoffset=-15mm%                  For Eng.


\setcounter{page}{15}


\god{2003}


\nomer{11~(498)} %                        Remove in Eng.


%\nomer{Vol.\,47, No.\,11}%               Set in Eng.


%\str{\pageref{begin}--\pageref{end}}%  Set in Eng.


\udk{514.763}





\author{\it Мирьяна Джорич\,${}^1$, Масафуми Окумура}


% NB:   ^^ remove in Eng.


\title{О кривизне $CR$-подмногообразий комплексного проективного


пространства, имеющих максимальную $CR$-размерность}


\thanks{${}^1$\,Работа выполнена автором при финансовой поддержке


Министерства науки Сербии


(проект ``{\it Геометрия, образование и визуализация с


приложениями}'').}





\date{}


%\pagestyle{myheadings}%         Set in Eng.


\pagestyle{plain} %              Remove in Eng.


\begin{document}


%\thispagestyle{plain}%          Set in Eng.


\maketitle


\label{begin}





\setcounter{section}{-1}








Пусть $M$ --- $n$-мерное $CR$-подмногообразие $CR$-размерности


$\frac{n-1}{2}$ в комплексном проективном пространстве, т.\,е.


$M$ --- вещественное подмногообразие комплексного проективного


пространства такое, что для любой точки $x$ в $M$ максимальное


голоморфное  подпространство касательного пространства к $M$ в $x$


имеет размерность $n-1$. Тогда $M$ необходимо имеет нечетную


размерность и существует единичное векторное поле $\xi$, нормальное


к $M$ со свойством $JT(M)\subset T(M)\oplus \Span\{\xi\}$. В данной статье


изучаются такие подмногообразия в предположении, что векторное поле


$\xi$ параллельно в нормальной связности, и находится достаточное


условие для того, чтобы такое подмногообразие было открытым


подмножеством в геодезической сфере.





\section{Введение}





Вещественные


гиперповерхности в комплексном проективном пространстве активно


изучаются в течение многих лет. Хотя после пространств


постоянной кривизны комплексное проективное пространство можно


считать самым простым, его геометрия налагает существенные


ограничения на геометрию гиперповерхностей. Например, в


комплексном проективном пространстве не существует ни


омбилических, ни эйнштейновых вещественных гиперповерхностей


\cite{CecilRyan}. Р.\,Такаги \cite{Takagi}, дав классифицикацию


однородных вещественных гиперповерхностей в комплексном


проективном пространстве (что само по себе было очень


существенным достижением), ввел целый ряд


гиперповерхностей. Эти гиперповерхности привлекли внимание


геометров, которые стали исследовать их свойства и находить


характеристики различных подклассов классификации Р.\,Такаги


(см., напр., \cite{CecilRyan},


\cite{Kim}--\cite{Oktrans}). При этом ряд подклассов был


охарактеризован с помощью свойств оператора второй основной


формы, тензора Риччи и других дифференциально-геометрических


объектов. Более детальное обсуждение данного вопроса и обширную


библиографию можно найти в \cite{NirWells}.








В данной статье мы продолжаем вышеупомянутые исследования и


находим новое достаточное условие того, что


$n$-мерное вещественное подмногообразие коразмерности $p$


в комплексном проективном пространстве, голоморфное касательное


подпространство которого имеет максимальную размерность, есть


открытое подмножество геодезической сферы. Известно, 


что вещественная гиперповерхность комплексного


проективного пространства обладает почти контактной метрической


структурой, индуцированной комплексной структурой. Более того, в


\S\,1 доказываем, что геодезическая гиперсфера в


комплексном проективном пространстве является сасакиевым


многообразием, находим ее тензор кривизны и замечаем, что он имеет


вид тензора кривизны сасакиевой пространственной формы


\cite{Blair}. Поэтому если тензор кривизны $R$


подмногообразия $M$ имеет  вид


\begin{align*}


R(X,Y)Z&=\alpha\{g(Y,Z)X-g(X,Z)Y\}+\\


&+\beta\{g(FY,Z)FX-g(FX,Z)FY-2g(FX,Y)FZ\}+ \\


&+\gamma\{u(Y)u(Z)X-u(X)u(Z)Y+u(X)g(Y,Z)U-u(Y)g(X,Z)U\},


\end{align*}


где $\alpha$, $\beta$, $\gamma$ --- вещественнозначные функции


на $M$, $F$~--- кососимметрический эндоморфизм пространства


$T(M)$, $U$ ---  касательное векторное поле и


$u$ --- 1-форма на $M$ такие, что $g(U,X)=u(X)$, то


будем говорить, что  $M$ имеет кривизну {\it типа сасакиевой


пространственной формы}. Цель статьи --- доказательство


следующей теоремы.








\begin{mainthm}


Пусть $M$ --- связное


$n$-мерное $(n>2p-1$, $p\ge 2)$ $CR$-подмногообразие


$CR$-размерности $\frac{n-1}{2}$ в комплексном проективном


пространстве $P^{\frac{n+p}{2}}(\Bbb C)$. Если нормальное


векторное поле $\xi$ параллельно в нормальной связности и


$M$ имеет кривизну типа сасакиевой пространственной формы, то


$M$ есть открытое подмножество в геодезической сфере. 


\end{mainthm}








\section{Гиперповерхности комплексного проективного пространства}








Пусть $\wt M(c)$ --- пространство постоянной голоморфной


секционной кривизны $4c$, вещественная размерность которого


равна $2m$,


$\wt \nabla$~--- его связность Леви-Чивита.


Для погруженного подмногообразия


$f: M^{2m-1}\rightarrow \wt M$,


связность Леви-Чивита $\nabla$ индуцированной метрики и оператор


$A$ второй основной формы погружения характеризуются следующими


соотношениями:


\begin{align}


\wt\nabla_XY&=\nabla_XY+g(AX,Y)\xi,\\


\wt\nabla_X\xi&=-AX,\notag


\end{align}


где $\xi$ --- локальное поле единичной нормали.





В дальнейшем для краткости будем опускать упоминание об


иммерсии $f$. Пусть $J:T\wt M\rightarrow T\wt M$


--- комплексная структура со свойствами $J^2=-I$,


$\wt \nabla J=0$, $g(JX,JY)=g(X,Y)$. Oпределим


``структурный вектор'', так называемый вектор Хопфа \cite{Berndt}:


$$


W=-J\xi.


$$


Далее определим кососимметричный $(1,1)$-тензор $\varphi$:


\begin{equation}


JX=\varphi X+g(X,W)\xi,


\end{equation}


где $\varphi X$ --- тангенциальная проекция $JX$. Теперь из


эрмитовости метрики $g$ следует


$$


\varphi^2=-X+g(X,W)W, \quad \varphi W=0.


$$


Заменив $\varphi^2=-I$ на $W^\perp=\{X\in T(M)\mid g(X,W)=0\}$,


получаем, что ранг $\varphi$ равен $2m-2$ и


$\ker \varphi=\Span\{W\}$.


Такой оператор  $\varphi$ определяет почти контактную


метрическую структуру \cite{Blair}, а  $W^\perp$ называется


голоморфным распределением.


Cтандартным образом выводим уравнения Гаусса и Кодацци:


\begin{gather}


R(X,Y)Z=g(AY,Z)AX-g(AX,Z)AY +c\{( g(Y,Z)X-g(X,Z)Y+ \notag \\


+ g(\varphi Y,Z)\varphi X


-g(\varphi X,Z)\varphi Y +2 g(X,\varphi Y)\varphi Z\}, \\


(\nabla_X A)Y- (\nabla_Y A)X=c( g(X,W)\varphi Y-


g(Y,W)\varphi X+2g(X,\varphi Y)W).\notag


\end{gather}


Используя (1.1), (1.2), $\wt \nabla J=0$, получаем


\begin{equation}


(\nabla _X\varphi)Y=g(Y,W)AX-g(AX,Y)W.


\end{equation}





Пусть $P^m(\Bbb C)$ --- комплексное проективное пространство.


Приведем список гиперповерхностей в $P^m(\Bbb C)$, которые


настолько важны для рассматриваемой теории, что вошли в


стандартный набор примеров. Эти гиперповерхности разделяются на


пять типов A--E, причем тип A разделяется еще на два подтипа:


\begin{itemize}


\item[(A1)] геодезическая сфера (трубчатая гиперповерхность вокруг


гиперплоскости $P^{m-1} (\Bbb C)$);


\item[(A2)] трубчатая поверхность вокруг вполне геодезического


подпространства $P^k(\Bbb C)$, $1\le k\le m-2$;


\item[(B)] трубчатая поверхность вокруг комплексной квадрики $Q^{m-1}$;


\item[(C)] трубчатая поверхность вокруг


$\,P^1(\Bbb C)\times P^k(\Bbb C)$,


при этом $2k+1=m$ и $m\ge 5$;


\item[(D)] трубчатая поверхность вокруг комплексного грассманова


многообразия $G_{2,5}$, при этом $m=9$;


\item[(E)] трубчатая поверхность вокруг эрмитова симметрического


пространства $SO(10)/U(5)$, при этом $m=15$.


\end{itemize}





Как было доказано в \cite{Takagi}, этот список состоит в


точности из однородных вещественных гиперповерхностей в


$\,P^m(\Bbb C)$ (его часто называют ``списком Такаги'').


Кроме того, каждая гиперповерхность Хопфа (т.\,е.


гиперповерхность, для которой вектор $W$ является главным) с


постоянными главными кривизнами есть открытое подмножество одной


из этих поверхностей. В доказательство этого результата внесли


вклад многие авторы, в окончательном виде он был получен


в \cite{Kimura}.





Теперь пусть $M$ --- геодезическая гиперсфера в комплексном


проективном пространстве $\wt M(c)=P^m(\Bbb C)$. Так как


геодезические сферы (тип A1) в комплексном проективном


пространстве имеют две различные главные кривизны


\cite{CecilRyan}:


$\frac{1}{r}\cot u$


кратности $2m-2$ и $\frac{2}{r}\cot 2u$ кратности $1$, то можем


записать $AX=\lambda X+\mu g(W,X)W$, и, используя (1.4), получим


$$


(\nabla _X\varphi)Y=\lambda\{g(Y,W)X-g(X,Y)W\},


$$


т.\,е. $M$ есть сасакиево многообразие \cite{Blair}. Более


того, используя уравнение Гаусса (1.3), получаем


\begin{align*}


R(X,Y)Z&=(c+\lambda^2)\{g(Y,Z)X-g(X,Z)Y\}+ \\ \notag


&+c\{g(\varphi Y,Z)\varphi X


-g(\varphi X,Z)\varphi Y


-2g(\varphi X,Y)\varphi Z\}+ \\


&+\lambda\mu\{g(Y,W)g(u(Z,W)X-g(X,W)g(Z,W)Y+ \\


&+g(X,W)g(Y,Z)W-g(Y,W)g(X,Z)W\},


\end{align*}


т.\,е. $R$ есть тензор кривизны сасакиевой пространственной


формы.








\section{$CR$-подмногообразия комплексной пространственной формы,\protect\\


имеющие максимальную $CR$-размерность}








В дальнейшем будем предполагать, что $\overline M$ ---


$(n+p)$-мерное келерово многообразие с келеровой структурой


$(J,\ov g)$ и  $M$ --- $n$--мерное вещественное подмногообразие


в $\overline M$, погруженное с помощью иммерсии $\imath$.


Тогда касательное расслоение $T(M)$ отождествляется с


подрасслоением в $T(\overline M)$ и риманова метрика $g$ на


$M$ индуцируется  римановой метрикой $\ov g$ на $\overline M$:


$g(X,Y)=\overline g(\imath X,\imath Y)$, где $X,Y\in T(M)$.


Дифференциал иммерсии обозначим также через $\imath$.


Нормальное расслоение $T^\perp(M)$ есть подрасслоение расслоения


$T(\overline M)$, состоящее из всех векторов


$\overline X\in T(\overline M)$, ортогональных


$T(M)$ в римановой метрике $\ov g$.





Подпространство $H_x(M)=JT_x(M) \cap T_x(M)$ есть максимальное


$J$-инвариантное подпространство касательного пространства


$T_x(M)$ в $x\in M$, называемое голоморфным касательным


пространством в $x$. Если размерность подпространства $H_x(M)$


не зависит от $x \in M$, то $M$ называется подмногообразием


Коши--Римана, или $CR$-подмногообразием, и комплексная размерность


$H_x(M)$ называется $CR$-размерностью подмногообразия $M$


\cite{NirWells}, \cite{Tumanov}. Известно, что вещественная


гиперповерхность является одним из типичных примеров


$CR$-подмногообразий,  $CR$-размерность которых равна


$\frac{n-1}{2}$, где $n$ --- размерность гиперповерхности.


Заметим, что в случае максимальной $CR$-размерности


вышеприведенное определение $CR$-подмногообразия совпадает с


определением $CR$-подмногообразия, данного в


\cite{Bejancu}. Подробное изложение с примерами


$CR$-подмногообразий максимальной $CR$-размерности см. в


\cite{Banach}.








Пусть  $M$ --- $CR$-подмногообразие максимальной $CR$-размерности.


Это значит, что в каждой точке $x \in M$ касательное пространство


$T_x(M)$ удовлетворяет условию


$\dim_{\Bbb R}(JT_x(M)\cap T_x(M))=n-1\,$.


Отсюда следует, что $M$ нечетномерно и существует единичное


векторное поле $\xi$, нормальное к $M$ и такое, что


$JT_x(M)\subset T_x(M)\oplus \Span\{\xi_x\}$ для любой $x\in M$.


Следовательно, для любого касательного векторного поля


$X$, выбрав локальный орторепер $\xi,\xi_1,\dotsc,\xi_{p-1}$


нормального пространства к $M$, имеем следующее разложение на


касательную и нормальную составляющую:


\begin{align}


J\imath X &=\imath\,FX+u(X)\xi, \\


J\xi &=-\imath\,U+P\xi, \\


J\xi_a &=-\imath\,U_a+P\xi_a\quad (a=1,\dotsc,p-1).


\end{align}


Здесь $F$ и $P$ являются кососимметричными эндоморфизмами


пространств $T(M)$ и $T^\perp(M)$ соответственно;


$U$, $U_a$, $a=1,\dotsc,p-1$, --- касательные векторные поля


и $u$ --- 1-форма на $M$. Используя (2.1)--(2.3), из


эрмитовости  $J$ получаем


\begin{gather*}


g(U,X)=u(X),\quad U_a=0\quad (a=1,\dotsc,p-1),\\


F^2X=-X+u(X)U, \\


u(FX)=0,\quad FU=0,\quad P\xi=0.


\end{gather*}


Следовательно, соотношения (2.2) и (2.3) можем записать в виде


\begin{equation}


J\xi=-\imath U,\quad J\xi_a=P\xi_a \quad


(a=1,\dotsc,p-1).


\end{equation}





Подпространство  $\{\eta\in T^\perp, \eta\perp\xi\}$ является


$J$-инвариантным, поэтому мы можем взять орторепер пространства


$T^\perp(M)$ вида


$\xi,\xi_1,\dotsc,\xi_q,\xi_{1^*},\dotsc,\xi_{q^*}$, где


$\xi_{a^*}=J\xi_a$ и $q=\frac{p-1}{2}$.


Обозначим через $\overline\nabla$ и


$\nabla$ римановы связности многообразий $\overline M$ и


$M$ соответственно, а через $D$ --- нормальную связность,


индуцированную связностью $\overline\nabla$ в нормальном


расслоении многообразия  $M$. Они связаны известным уравнением


Гаусса


\begin{equation}


\overline\nabla_{\imath X}\imath Y=\imath


\nabla_XY+h(X,Y),


\end{equation}


где через $h$ обозначена вторая основная форма, и уравнениями


Вейнгартена


\begin{align*}


\overline \nabla_{\imath X}\xi &=


-\imath AX + D_X\xi -\imath AX +


\sum_{a=1}^q\{s_a(X)\xi_a + s_{a^*}(X)\xi_{a^*}\}, \\


\overline \nabla_{\imath X}\xi_a &= -\imath A_aX + D_X\xi_a


= -\imath A_aX - s_a(X)\xi


+ \sum_{b=1}^q\{ s_{ab}(X)\xi_b + s_{ab^*}(X)\xi_{b^*}\},\\


\overline \nabla_{\imath X}\xi_{a^*} &= -\imath A_{a^*}X + D_X\xi_{a^*}


= -\imath A_{a^*}X - s_{a^*}(X)\xi


+ \sum_{b=1}^q\{ s_{a^*b}(X)\xi_b + s_{a^*b^*}(X)\xi_{b^*} \},


\end{align*}


где  $A$, $A_a$, $A_{a^*}$ являются операторами второй основной


формы для нормалей


$\xi$, $\xi_a$, $\xi_{a^*}$ соответственно,


и через $s$ обозначены коэффициенты нормальной связности $D$.


Так как объемлющее многообразие келерово, имеем


\begin{gather}


A_{a^*}X = FA_aX - s_a(X)U, \\


s_{a^*}(X) = u(A_aX) = g(A_aX,U) = g(A_aU,X), \\


s_{a^*b^*} = s_{ab}, \quad s_{a^*b} = -s_{ab^*},\notag \\


h(X,Y) = g(AX,Y)\xi + \sum_{a=1}^q\{ g(A_aX,Y)\xi_a +


g(A_{a^*}X,Y)\xi_{a^*}\}


\end{gather}


для всех $X,Y\in T(M)$.





Далее ковариантно продифференцируем соотношения (2.1) и (2.2),


используя то, что почти комплексная структура $J$ инвариантна


относительно римановой связности $\overline\nabla$ на


$\overline M$, и, сравнив тангенциальные и нормальные


составляющие, с учетом (2.4) получим


\begin{gather*}


(\nabla_YF)X=u(X)AY-g(AY,X)U, \\


(\nabla_Yu)(X)=g(FAY,X), \\


\nabla_XU=FAX.


\end{gather*}





Если предположить, что векторное поле


$\xi$ параллельно относительно нормальной связности


$D$, то 


$$


D_X\xi=\sum_{a=1}^q\{s_a(X)\xi_a+s_{a^*}(X)\xi_{a^*}\}=0,


$$


откуда следует $s_a=s_{a^*}=0$ ($a=1,\dotsc,q$).


Теперь, используя соотношения (2.6) и (2.7), получаем


\begin{alignat}{2}


A_{a^*}&=FA_a  &\quad (a&=1,\dotsc,q), \\


A_aU&=0  &\quad (a&=1,\dotsc,q).


\end{alignat}





Так как вторая основная форма


$h(X,Y)$ симметрична относительно $X,Y$, из (2.8) и (2.9)


следует, что $FA_{a^*}$, $a=1,\dotsc,q$, симметричны, откуда


\begin{equation}


FA_a+A_aF=0 \quad (a=1,\dotsc,q).


\end{equation}





Наконец, если объемлющее многообразие $\overline M$ является


комплексной пространственной формой, т.\,е. келеровым


многообразием постоянной голоморфной секционной кривизны $4k$, то


$$


\overline R(\overline X,\overline Y)\overline Z=


k\{\ov g(\overline Y,\overline Z)\overline X-


\ov g(\overline X,\overline Z)\overline Y+


\ov g(J\overline Y,\overline Z)J\overline X-


\ov g(J\overline X,\overline Z)J\overline Y


-2\ov g(J\overline X,\overline Y)J\overline Z\}


$$


для  $\overline X$, $\overline Y$, $\overline Z$, касательных к


$\overline M$, где $\overline R$ --- тензор кривизны риманова


многообразия $\overline M$. Поэтому уравнения Гаусса и Кодацци,


записанные для единичного вектора $\xi$, имеют вид


\begin{gather}


R(X,Y)Z=k\{g(Y,Z)X-g(X,Z)Y+g(FY,Z)FX-g(FX,Z)FY+ \\


-2g(FX,Y)FZ\}+g(AY,Z)AX-g(AX,Z)AY+\notag \\


+\sum_{a=1}^q\{g(A_aY,Z)A_aX-g(A_aX,Z)A_aY+\notag \\


+g(FA_aY,Z)FA_aX-g(FA_aX,Z)FA_aY\},\notag\\


(\nabla_YA)X-(\nabla_XA)Y=k\{u(Y)FX-u(X)FY+


2g(FX,Y)U\}\notag


\end{gather}


соответственно для всех $X$, $Y$, $Z$, касательных к $M$, где через


$R$ обозначен  тензор кривизны риманова многообразия $M$


(\cite{KobNom}).








\section{Оператор второй основной формы $CR$-подмногообразия\protect\\


максимальной $CR$-размерности}








В начале этого параграфа напомним некоторые результаты.


Пусть $M$ --- $n$-мерное ($n\ge 3$)


$CR$-подмногообразие $CR$-размерности $\frac{n-1}{2}$ в


$(n+p)$-мерном келеровом многообразии  $\overline M$ постоянной


голоморфной секционной кривизны $4k$. Предположим, что оснащающее


векторное поле $\xi$ подмногообразия $M$ параллельно относительно


нормальной связности $D$. Тогда, если оператор второй основной


формы $A$ относительно $\xi$ имеет единственное собственное


значение, то $k=0$ (\cite{DjoricOkumura}, лемма 3.3). Более того,


если  $k\ne 0$ и оператор второй основной формы


$A$ имеет в точности два различных собственных значения, то $U$


есть собственный вектор оператора $A$ (\cite{DjoricOkumura}, лемма 3.4).


Далее, пусть $\overline M$ ---


комплексное проективное пространство


$P^{\frac{n+p}{2}}(\Bbb C)$ постоянной голоморфной секционной кривизны


$4k$, где $k>0$ и $n>3$. Тогда если оператор второй основной формы


$A$ имеет в точности два различных собственных значения, то


эти значения постоянны (\cite{DjoricOkumura}, лемма 4.1). Если


дополнительно предположить, что $n>2p-1$, $p\ge 2$, то кратность


собственного значения, соответствующего собственному вектору $U$


оператора $A$, равняется единице.





Пусть объемлющее многообразие $\overline M$ есть


комплексное проективное пространство


$P^{\frac{n+p}{2}}(\Bbb C)$, наделенное метрикой Фубини--Штуди


постоянной голоморфной секционной кривизны $4$. Пусть


$\Gamma(x,\xi,r)$, $-\infty <r<\infty$, --- геодезическая в


$P^{\frac{n+p}{2}}(\Bbb C)$, параметризованная длиной дуги $r$,


причем $\Gamma(x,\xi,0)=x\in P^{\frac{n+p}{2}}(\Bbb C)$ и


$\frac{d}{dr}\Gamma(x,\xi,0)=\xi$. Определим отображение $\Phi_r$ из $M$


в $P^{\frac{n+p}{2}}(\Bbb C)$:


$$


\Phi_r(x)=\Gamma(x,\xi,r).


$$


Можно найти  $r=r_0$ такое, что


$(\Phi_{r_0})_*X=0$ для всех $X\in T(M)$


\cite{DjoricOkumura}. Следовательно, отображение $\Phi_{r_0}$


постоянно, т.\,е. $\Phi_{r_0}(M)=y\in P^{\frac{n+p}{2}}(\Bbb C)$.


Так как $r_0$ есть длина дуги геодезической, соединяющей


$x$ и $y$, то любая точка $x\in M$ лежит на геодезической


гиперсфере с центром  $y$ и радиусом $r_0$. Отсюда следует





\begin{athm}[{\series{m}\shape{n}\selectfont\cite{DjoricOkumura}}]


Пусть $M$ --- $n$-мерное $(n>2p-1$, $p\ge 2)$ $CR$-подмногообразие


$CR$-размерности $\frac{n-1}{2}$ в комплексном проективном


пространстве $P^{\frac{n+p}{2}}(\Bbb C)$. Если оператор второй


основной формы $A$ относительно оснащающего нормального


векторного поля $\xi$ имеет в точности два собственных


значения и $\xi$ параллельно в нормальной связности, то


существует геодезическая сфера


$S$ в $P^{\frac{n+p}{2}}(\Bbb C)$, на которой лежит $M$.


\end{athm}





\begin{lem}


Пусть  $M$ --- связное


$n$-мерное $(n>3)$  $CR$-подмногообразие $CR$-размерности


$\frac{n-1}{2}$ в комплексном проективном пространстве


$P^{\frac{n+p}{2}}(\Bbb C)$.


Если оснащающее нормальное векторное поле


$\xi$ параллельно в нормальной связности и тензор


кривизны $R$ многообразия $M$ имеет вид


\begin{align}


R(X,Y)Z&=\alpha\{g(Y,Z)X-g(X,Z)Y\}+\notag \\


&+\beta\{g(FY,Z)FX-g(FX,Z)FY-2g(FX,Y)FZ\}+\notag \\


&+\gamma\{u(Y)u(Z)X-u(X)u(Z)Y+u(X)g(Y,Z)U-u(Y)g(X,Z)U\},


\end{align}


где  $\alpha$, $\beta$, $\gamma$ --- вещественнозначные функции


на $M$ $($т.\,е. $M$ имеет кривизну типа сасакиевой


пространственной формы$)$, то оператор второй основной формы


$A$ относительно $\xi$ имеет в точности два собственных значения


и $U$ есть собственный вектор оператора $A$. 


\end{lem}





\begin{pf}


Вначале подсчитаем $R(X,Y)Z$ двумя


способами: используя условие (3.1) и уравнение Гаусса


(2.12). Положив $Z=U$ и используя (2.10), получаем


\begin{equation}


(\alpha+\gamma-1)(u(Y)X-u(X)Y)=g(AY,U)AX-g(AX,U)AY


\end{equation}


для любых $X,Y\in T(M)$. Так как оператор второй основной формы


$A$ относительно $\xi$ имеет по крайней мере два различных


собственных значения, например $\lambda_1$ и $\lambda_2$,


то, используя соотношение (3.2), получаем


$$


\lambda_1(x)\lambda_2(x)=\alpha(x)+\gamma(x)-1, \quad x\in M.


$$


Поскольку $M$ связно, то оператор


$A$ имеет в точности два собственных значения, которые постоянны, а


$U$ --- собственный вектор оператора $A$.


\end{pf}








\begin{notenonum}


В предположениях леммы 1, если


$n>2p-1$, $p\ge 2$, то кратность собственного значения


$\mu$, соответствующего собственному вектору $U$ оператора второй


основной формы $A$, равна единице  и


\begin{equation}


AX=\lambda X+(\mu -\lambda)u(X)U.


\end{equation}


\end{notenonum}





Из леммы 1 и теоремы A получается








\begin{thm}


Пусть $M$ --- связное


$n$-мерное $(n>2p-1$, $p\ge 2)$ $CR$-под\-мно\-го\-об\-разие


$CR$-размерности $\frac{n-1}{2}$ в комплексном проективном


пространстве $P^{\frac{n+p}{2}}(\Bbb C)$. Если оснащающее


нормальное векторное поле $\xi$ параллельно в нормальной


связности и $M$ имеет кривизну типа сасакиевой пространственной


формы, то существует геодезическая сфера


$S$ в пространстве


$P^{\frac{n+p}{2}}(\Bbb C)$, на которой лежит многообразие  $M$.


\end{thm}








\begin{lem}


Пусть $M$ --- $n$-мерное $(n>2p-1$,


$p\ge 2)$ $CR$-подмногообразие $CR$-размерности


$\frac{n-1}{2}$ в комплексном проективном пространстве


$P^{\frac{n+p}{2}}(\Bbb C)$.


Если оснащающее векторное поле $\xi$ параллельно в нормальной


связности и  $M$ имеет кривизну типа сасакивой пространственной


формы, то


$A_a=A_{a^*}=0$, $a=1,\dotsc,q$, где $A_a$, $A_{a^*}$


--- операторы второй основной формы относительно нормалей


$\xi_a$, $\xi_{a^*}$ соответственно. 


\end{lem}








\begin{pf}


Для доказательства леммы используем


результаты \cite{ONeill}. Определим симметрическую


билинейную форму


\begin{gather*}


k: T_0(M)\times T_0(M)\to


\Span\{\xi_1,\dotsc,\xi_q,\xi_{1^*},\dotsc,\xi_{q^*}\}, \\


k(X,Y)=h(X,Y)-g(AX,Y)\xi=


\sum_{a=1}^q\{g(A_aX,Y)\xi_a+g(A_{a^*}X,Y)\xi_{a^*}\},


\end{gather*}


где $T_0(M)=\{X\in T(M)\mid g(X,U)=0\}$. Заметим, что на $T_0(M)$


оператор $F$ является почти комплексной структурой, т.\,е.


$F^2=-I$ на $T_0(M)$.





Так как подрасслоение


$T_1^\perp(M)=\{\eta\in T^\perp(M)\mid\overline g(\eta,\xi)=0\}$


нормального расслоения $T^\perp(M)$ является


$J$-инвариантным, используя соотношения  (2.1), (2.4) и уравнение


(2.5), получим


\begin{equation}


Jk(X,Y)=k(X,FY),


\end{equation}


откуда следует, что билинейная форма $k$ является почти


комплексной. Дискриминант $\Delta$ формы $k$ на плоскости


$\Span\{X,Y\}$ имеет вид


$$


\Delta_{X Y}=\frac{\overline g(k(X,X),k(Y,Y))-\|k(X,Y)\|^2}


{g(X,X)g(Y,Y)-g(X,Y)^2}.


$$


Тогда в силу (3.4) для единичного векторного поля


$X$, принадлежащего $T_0(M)$, имеем


$$


\Delta_{X FX}=-2\|k(X,Y)\|^2.


$$





Заметим, что $\Delta_{X FX}$ в точности соответствует


голоморфной разности $\Delta_{\mathrm{hol}}$, определенной  в


\cite{ONeill}. Далее, используя уравнение Гаусса


(2.12) и соотношение (2.11), имеем


$$


g(R(X,FX)FX,X)=4+g(AFX,FX)g(AX,X)-g(AFX,X)^2-2\|k(X,X)\|^2.


$$


Более того, используя соотношение (3.3), получаем


\begin{equation}


g(R(X,FX)FX,X)=4+\lambda^2-2\|k(X,X)\|^2.


\end{equation}


Теперь, используя  условие (3.1), имеем


\begin{equation}


g(R(X,FX)FX,X)=\alpha+3\beta.


\end{equation}


Поэтому, сравнивая соотношения (3.5) и (3.6), получаем


$$


\|k(X,X)\|^2=\dfrac{1}{2}(4+\lambda^2-\alpha-3\beta),


$$


т.\,е. дискриминант $\Delta_{X FX}$ постоянен, и из леммы 6


\cite{ONeill} следует, что $k$ изотропна. Наконец, используя


лемму 8 из \cite{ONeill}, получаем, что


$\frac{n-1}{2}\big(\frac{n-1}{2}+1\big)=\frac{n^2-1}{4}$


векторов $k(e_i,e_j)$, $Jk(e_i,e_j)$ ортогональны, где


$e_1,\dotsc,e_{\frac{n-1}{2}},Je_1,\dotsc,Je_{\frac{n-1}{2}}$ ---


орторепер пространства $T_0(M)$. Однако из


$p<\frac{n+1}{2}$ следует


$\frac{n^2-1}{4}>p-1$, поэтому $k(e_i,e_j)=0$, т.\,е.


$h(X,Y)=g(AX,Y)\xi$,


что эквивалентно равенствам $A_a=0$,


$A_{a^*}=0$, $a=1,\dotsc,q$.


\end{pf}








\section{Редукция коразмерности и характеристический признак


геодезической гиперсферы}








Пусть $M$ --- $n$-мерное ($n>2p-1$,


$p\ge 2$) $CR$-подмногообразие $CR$-размерности


$\frac{n-1}{2}$ в комплексном проективном пространстве


$P^{\frac{n+p}{2}}(\Bbb C)$,


оснащающее векторное поле $\xi$ параллельно в нормальной


связности, и $M$ имеет кривизну типа сасакивой пространственной


формы.


Пусть $N_0(x)=\{\xi\in T^\perp_x(M)\mid A_{\xi}=0\}$ и


$H_0(x)$ --- максимальное $J$-инвариантное подпространство в $N_0(x)$,


т.\,е. $H_0(x)=JN_0(x)\cap N_0(x)$. Нам потребуется








\begin{athm}[{\series{m}\shape{n}\selectfont\cite{Okumura}}]


Пусть $M$ --- $n$-мерное вещественное подмногообразие комплексного 


проективного пространства $P^{\frac{n+p}{2}}(\Bbb C)$ вещественной


размерности $(n+p)$. Если ортогональное дополнение


$H_1(x)$ подпространства  $H_0(x)$ в $T^\perp(M)$


инвариантно относительно параллельного переноса в нормальной


связности и $r$ --- постоянная размерность пространства


$H_1(x)$, то существует  вполне геодезическое


комплексное проективное пространство


$P^{\frac{n+r}{2}}(\Bbb C)$ вещественной размерности $(n+r)$


такое, что $M\subset P^{\frac{n+r}{2}}(\Bbb C)$. 


\end{athm}








При вышеуказанных условиях на $M$ из леммы 2 следует


$A_a=A_{a^*}=0$ при $a=1,\dotsc,q$, поэтому получаем


$\Span\{\xi_1(x), \dotsc, \xi_q(x),$  $\xi_{1^*}(x), \dotsc,


\xi_{q^*}(x)\} \subset N_0(x)$.


С другой стороны, для любой $\eta\in N_0(x)$ положим


$\eta=p_0\xi+\sum\limits_{a=1}^q\{p_a \xi_a+p_{a^*}\xi_{a^*}\}$.


Тогда $A_{\eta}=p_0A + \sum\limits_{a=1}^{q} \{p_a A_a+p_{a^*}A_{a^*}\}


=p_0A=0$. Отсюда $p_0=0$ и


$$


\eta=\sum_{a=1}^{q}\{p_a\xi_a+p_{a^*}\xi_{a^*}\}\in


\Span \{\xi_1(x),\dotsc,\xi_{q}(x),


\xi_{1^*}(x),\dotsc,\xi_{q^*}(x)\}.


$$


Поэтому $N_0(x)=\Span


\{\xi_1(x),\dotsc,\xi_{q}(x),\xi_{1^*}(x),\dotsc,\xi_{q^*}(x)\}$.





Далее, используя второе уравнение в (2.4), получаем


$JN_0(x)=N_0(x)$, и, следовательно,


$H_0(x)=\Span\{\xi_1(x),\dotsc,\xi_{q}(x),\xi_{1^*}(x),


\dotsc,\xi_{q^*}(x)\}$.


Поэтому ортогональное дополнение $H_1(x)$ подпространства


$H_0(x)$ в $T^\perp(M)$ есть $\Span\{\xi\}$.


Тогда из вышеуказанных условий на $M$ следует, что


$H_1(x)$ инвариантно относительно параллельного переноса


в нормальной связности, и из теоремы~B о редукции коразмерности


получаем следующее утверждение.








\begin{thm}


Пусть $M$ ---


$n$-мерное $(n>2p-1$, $p\ge 2)$ $CR$-подмногообразие $CR$-размерности


$\frac{n-1}{2}$ в комплексном проективном пространстве


$P^{\frac{n+p}{2}}(\Bbb C)$.


Если оснащающее нормальное векторное поле


$\xi$ параллельно относительно нормальной связности


и $M$ имеет кривизну типа сасакиевой пространственной формы, то


существует вполне геодезическое комплексное проективное


пространство $P^{\frac{n+1}{2}}(\Bbb C)$ вещественной


размерности $n+1$ такое, что


$M\subset P^{\frac{n+1}{2}}(\Bbb C)$.


\end{thm}








Теперь, используя теорему 2, можно рассматривать подмногообразие


$M$ как вещественную гиперповерхность в


$P^{\frac{n+1}{2}}(\Bbb C)$, являющуюся вполне геодезическим


подмногообразием  в $P^{\frac{n+p}{2}}(\Bbb C)$. 


Обозначим $P^{\frac{n+1}{2}}(\Bbb C)$ через $M^\prime$,


погружение $M$ в $M^\prime$ --- через $\imath_1$, 


вполне геодезическое погружение


$M^\prime$ в $P^{\frac{n+p}{2}}(\Bbb C)$ --- через $\imath_2$.


Тогда из (2.5) следует


$$


\nabla^\prime_{\imath_1 X}\imath_1Y=


\imath_1 \nabla_XY+h^\prime(X,Y)=\imath_1\nabla_XY+


g(A^\prime X,Y)\xi^\prime,


$$


где $h^\prime$~--- вторая основная форма $M$ в $M^\prime$,


$A^\prime$ --- соответствующий оператор второй основной формы


и $\xi^\prime$ --- единичное векторное поле на $M^\prime$,


нормальное к $M$.





Так как $\imath=\imath_2\cdot \imath_1$, имеем


\begin{equation}


\ov \nabla_{\imath_2\cdot \imath_1 X}\imath_2\cdot


\imath_1Y=\imath_2 \nabla^\prime_{\imath_1 X}\imath_1Y+


\ov h(\imath_1 X,\imath_1 Y)=\imath_2(\imath_1


\nabla_XY+g(A^\prime X,Y)\xi^\prime)


\end{equation}


в силу того, что  $M^\prime$ является вполне геодезическим в


$P^{\frac{n+p}{2}}(\Bbb C)$. Сравнивая соотношение (4.1) с


(2.5), получаем $\xi=\imath_2 \xi^\prime$, $A=A^\prime$.


Более того, т.\,к. $M^\prime$ есть комплексное подмногообразие в


$P^{\frac{n+p}{2}}(\Bbb C)$, для любого $X^\prime\in T(M^\prime)$


выполняется равенство


$J\imath_2X^\prime=\imath_2J^\prime X^\prime$,


где $J^\prime$ --- индуцированная комплексная структура на


$M^\prime=P^{\frac{n+1}{2}}(\Bbb C)$. Таким образом, из


(2.1) следует


$$


J\imath X=J\imath_2\cdot \imath_1 X=\imath_2J^\prime \imath_1


X=\imath_2(\imath_1 F^\prime X+\nu^\prime(X)\xi^\prime) 


=\imath F^\prime X+\nu^\prime(X)\imath_2\xi^\prime=\imath


F^\prime X+V^\prime\xi.


$$


Сравнивая это уравнение с (2.1), получаем


$$


F=F^\prime,\quad V^\prime=U, \quad \nu^\prime=u.


$$





Далее, в силу теоремы 2 \ $M$ есть вещественная гиперповерхность


в $P^{\frac{n+1}{2}}(\Bbb C)$, у которой оператор второй основной


формы $A$ относительно оснащающего векторного поля $\xi$ имеет в


точности два собственных значения. Теперь основная теорема


вытекает из следующей теоремы.








\begin{athm}[{\series{m}\shape{n}\selectfont\cite{CecilRyan}}]


Пусть $M$ --- связная вещественная


гиперповерхность в $\Bbb CP^m$,


$m\ge 3$, имеющая не более двух различных главных кривизн в


каждой точке. Тогда $M$ есть открытое подмножество геодезической


гиперсферы. 


\end{athm}








\begin{thebibliography}{99}





\bibitem{CecilRyan} %1


Cecil T.E., Ryan P.J.


{\em Focal sets and real


hypersurfaces in complex projective space} // Trans. Amer. Math.


Soc. -- 1982. -- V.\,269. -- \No\,2. -- P.\,481--499.





\bibitem{Takagi} %2


Takagi R.


{\em  On homogeneous real hypersurfaces in a complex projective


space} // Osaka J. Math. -- 1973. -- V.\,10. -- \No\,3. -- P.\,495--506.





\bibitem{Kim} %3


Kimura M.


{\em   Sectional curvatures of holomorphic planes on a real


hypersurface in $P^n(C)$} // Math. Ann. -- 1987. -- V.\,276. -- P.\,487--497.





\bibitem{Lawson} %4


Lawson H.B.


{\em Local rigidity theorem in rank-$1$ symmetric spaces} //


J. Different. Geom. -- 1970. -- V.\,4. -- P.\,349--357.





\bibitem{Oktrans} %5


Okumura M.


{\em On some real hypersurfaces of a complex projective space} //


Trans. Amer. Math. Soc. -- 1975. -- V.\,212. -- P.\,355--364.





\bibitem{NirWells} %6


Nirenberg R., Wells R.O. Jr.


{\em  Approximation theorems on differentiable submanifolds of a


complex manifold} // Trans. Amer. Math. Soc. -- 1965. -- V.\,142.


-- P.\,15--35.





\bibitem{Blair} % 7


Blair D.E.


{\em Contact manifolds in Riemannian geometry} //


Lect. Notes Math. -- Berlin: Springer, 1976. -- V.\,509. -- 146 p.





\bibitem{Berndt} %8


Berndt J.


{\em Uber Untermanningfaltigkeiten von komplexen Raumformen}:


Doct. Dissert. -- Univ. zu K\"oln, 1989.





\bibitem{Kimura} %9


Kimura M.


{\em   Real hypersurfaces


and complex submanifolds in complex projective space} // Trans.


Amer. Math. Soc. -- 1986. -- V.\,296. -- \No\,1. -- P.\,137--149.





\bibitem{Tumanov} %10


Tumanov A.E.


{\em  The geometry of CR manifolds}.


Encyclopedia of Math. Sci.  9.\,VI, Several complex variables III. --


Springer-Verlag, 1986. -- P.\,201--221.








\bibitem{Bejancu} %11


Bejancu A. {\em CR-submanifolds of a Kaehler manifold.} I //


Proc. Amer. Math. Soc. -- 1978. -- V.\,69. -- \No\,1. -- P.\,135--142.





\bibitem{Banach} %12


Djori\'c M., Okumura M.


{\em CR submanifolds of maximal CR dimension in complex manifolds} //


PDE's, Submanifolds and Affine Different. Geom., Banach


Center Publ., Inst. Math., Polish Acad.


Sci. -- Warszawa, 2002. -- V.\,57. -- P.\,89--99.





\bibitem{KobNom} %13


Kobayashi S., Nomizu K.


{\em Foundations of differential geometry.} II. --


New~York: Intersci., 1969. -- 470 p.





\bibitem{DjoricOkumura} %14


Djori\'c M., Okumura M.


{\em CR submanifolds of


maximal CR dimension of complex projective space} // Arch. Math. -- 1998. --


V.\,71. -- P.\,148--158.








\bibitem{ONeill} %15


O'Neill B.


{\em  Isotropic and K\"ahler immersions} //


Canad. J. Math. -- 1965. -- V.\,17. -- P.\,907--915.








\bibitem{Okumura} %16


Okumura M.


{\em Codimension reduction problem for real submanifolds of complex


projective space} // Colloq. Math. Soc. J\'anos Bolyai. -- 1989. -- V.\,56.


-- P.\,574--585.





\end{thebibliography}





\vskip 0.5cm





\post{\it Белградский университет $($Сербия$)$}{\it Поступила}


\post{\it Университет Сайтама $($Япония$)$}{15.01.2003}





\label{end}


\end{document}














%\def \Ime#1#2{\noindent #1\hfill #2\break}


%\Ime{Faculty of Mathematics} {Department of Mathematics}


%\Ime{University of Belgrade}{Saitama University}


%\Ime{Studentski trg 16, pb 550}{Shimo-Okubo}


%\Ime{11000 Belgrade, Yugoslavia}{Urawa, 338 Japan}


%\Ime{mdjoric$\@$matf.bg.ac.yu}{mokumura$\@$h8.dion.ne.jp}

















\bibitem{Debrecen}


 M.~Djori\'c and M.~Okumura


  {\em Certain application of an integral formula to


CR submanifold of complex projective space} to appear in Publ.


Math. Debrecen, V.62, 2003.





\bibitem{Klingenberg}


W.~Klingenberg


 {\em Real hypersurfaces in K\"ahler manifolds}


SFB 288, Preprint No.261. Diferentialgeometrie und Quantenphysik





\bibitem{NiebRyan}


R.~Niebergall and P.J.~Ryan


{\em  Real hypersurfaces in complex space forms}, {\rm in}


Tight and taut submanifolds


 (eds. T.E. Cecil and S.-s. Chern), Math. Sciences Res. Inst.


Publ. 32, Cambridge Univ. Press, Cambridge,  1997, 233--305.





